\documentclass[11pt]{article}
\usepackage[margin=1in]{geometry}
\usepackage[T1]{fontenc}
\usepackage[utf8]{inputenc}
\usepackage{url}
\setlength{\parindent}{0pt}
\setlength{\parskip}{0.7em}
\begin{document}

\today

Dear Dr. Steven K. Morley,
Editor-in-Chief, Space Weather
American Geophysical Union

Re: Submission of our manuscript ``Scale-dependent WSPR response to
geomagnetic storms: frequency dependence and network-size effects on
absolute loss'' for consideration in \emph{Space Weather}.

We are an independent family research group (Alexey Petrov and Makar
Petrov, no institutional affiliation). We present an 18-month,
three-band (20/40/15 m) quantitative analysis of WSPR spot counts
during geomagnetic storms ($Kp \ge 5$), using openly available
citizen-science data from wspr.live together with GFZ Potsdam $Kp$
and GOES X-ray indices. We believe this work aligns with your vision
of Space Weather as a maturing applied field, bridging open
citizen-science data with rigorous statistical controls.

Main findings:

\begin{enumerate}
\item Storm-time suppression of spots per active transmitter on all
three bands ($-23.9\%$, $-22.5\%$, $-30.1\%$; block-permutation
$p \le 2\times10^{-5}$), with an independent held-out replication
(Jan--Mar 2026).
\item A statistically robust auroral-to-equatorial gradient of
absorption: $-72.7\%$ at high latitudes versus $-12.8\%$ at low
latitudes, a $5.7\times$ ratio.
\item A selection-bias control: normalising hourly counts by the
number of active transmitters reveals that the long-standing
``15 m low-latitude anomaly'' is an artifact of operator behaviour
rather than a propagation effect.
\item Two distinct loss channels: MUF suppression and D-layer
absorption, joined by a decode-channel mechanism --- storm-time
Doppler spreading of the narrow ($\sim$1.4 Hz) WSPR tones, which
removes spots even when received SNR is high.
\end{enumerate}

The analysis has passed a strict independent audit: a multi-round
Q1 narrative and consistency review (estimate families separated,
effective sample size recomputed via $\tau_{\rm int}$, permutation
floors corrected), and a community critique by Gwyn Griffiths
(G3ZIL) of the HamSCI community --- 15 comments, all addressed; a
point-by-point response is attached as a supplementary file.
Statistical methodology includes block permutation with IAAFT
surrogates, Benjamini--Hochberg FDR control, and episode-level
robustness checks.

All data and code are openly available: the archived data package
(Zenodo v11) is at \url{https://doi.org/10.5281/zenodo.23802881}
(concept DOI: \url{https://doi.org/10.5281/zenodo.23285958}); the
analysis code is at \url{https://github.com/FelixRLEPERS/crosscorr}
(MIT license).

This manuscript is original, has not been published elsewhere, and
is not under consideration by another journal. All authors have
approved the submission. We declare no conflicts of interest.

We believe this bias-control method --- normalizing by active
transmitters --- is directly applicable to other citizen-science
datasets and may prevent false physical interpretations in future
WSPR studies.

We believe the results are of direct interest to the \emph{Space
Weather} readership: they quantify a practical HF-propagation impact
of geomagnetic storms using scalable, open citizen-science data, and
they demonstrate a bias-control method that the WSPR community can
apply directly.

Thank you for your consideration.

Sincerely,

Alexey Petrov (corresponding author) and Makar Petrov
Independent Researchers, Russia
felixrpetrov@gmail.com
https://github.com/FelixRLEPERS/crosscorr

\end{document}
